\chapter{相关技术介绍}

\section{前端开发技术与框架}

本系统前端采用 Vue 3 框架，并结合 Vite 构建工具进行开发。Vue 3 引入了 Composition API，使得逻辑复用和代码组织更加高效和灵活。系统在状态管理方面使用了 Pinia，路由管理采用 Vue Router 4。UI 组件库方面，系统融合了 Vuetify 3 与 Element Plus，利用其丰富的组件快速构建现代化的用户界面。在多媒体处理方面，系统利用 HTML5 Video 标签进行视频播放，并结合自定义事件实现与 AI 聊天的联动。

\section{后端开发技术与框架}

后端服务采用 Python 的 Flask 框架构建，Flask 的轻量级和可扩展性使其非常适合构建提供 RESTful API 的微服务。系统后端使用了 SQLAlchemy 作为 ORM 库进行关系型数据库的操作。在并发处理和后台任务方面，系统集成了线程池技术处理视频解析、文档 OCR 识别及向量化等长耗时任务。对于实时问答和生成式任务，系统利用 Flask 的流式响应特性（Server-Sent Events）实现与前端的打字机效果交互。

\section{大语言模型与检索增强生成（RAG）}

大语言模型（LLM）具有强大的自然语言理解与生成能力，但在特定领域或私有知识库应用中，容易产生“幻觉”。为了解决这一问题，系统引入了检索增强生成（RAG）技术。RAG 技术首先将本地知识库（如课程字幕、讲义文档）切片并转化为稠密向量，存入 FAISS 等向量数据库中；在用户提问时，系统先利用混合检索策略（词频统计的 BM25 加上语义向量检索）召回最相关的上下文片段，随后将其与用户问题拼接成提示词（Prompt）送入大模型，从而生成准确且有依据的回答。

\section{代码静态分析与启发式推理}

在智能代码实训模块中，系统并非简单依赖大模型直接改错，而是结合了代码静态分析（Static Analysis）技术。系统通过 Python 内置的 \texttt{ast} 模块，对学习者提交的代码进行语法解析，提取命名规范、圈复杂度、异常处理结构等质量指标，生成静态分析报告。随后，大模型基于预设的“苏格拉底式（Socratic）”系统提示词，结合该分析报告与题目信息，生成针对逻辑漏洞的反问和边缘测试用例引导。这种结合确定性分析与概率性推理的混合架构，是实现个性化启发式教学的核心技术。
